\section{Capabilities}
\label{sec:capabilities}

The previous two sections described structure. This one describes what the
structure enables, organized by the requirements of \secref{sec:motivation}.

\subsection{R2 in practice: late-bound live targets}

The single capability that most distinguishes \sys from a versioned prompt file is
that client code resolves \emph{which} version to run at call time. An agent loads
\prodalias{}; the platform decides what \prodalias{} points to. A remediation
therefore requires no client deployment, and the outer lane of
\figref{fig:system} exists to make exactly this visible.

The consequences are worth separating, because they are usually collapsed into a
vague claim about agility:

\begin{itemize}
\item \textbf{Remediation latency decouples from deployment latency.} The bound on
      time-to-fix becomes the refinement path plus operator attention, not the CI
      pipeline. We specify how to measure this in \appref{app:protocol} and have
      not measured it.
\item \textbf{Rollback is a pointer move, not a redeploy.} Where the governed path
      committed a checkpoint, restoring the recorded target is a bounded database
      and registry operation rather than a code revert. Where an authoring route
      moved the alias instead, there is no checkpoint to restore from
      (\secref{sec:limitations}).
\item \textbf{Authority is scope-checked.} Rotating an alias requires the
      \operator{} or \admin{} scope. We are careful not to call this separation of
      duties: the sibling \approver{} scope exists, but nothing currently requires
      that the principal who applies a change differ from the one who authored it,
      and the approval record is minted already-approved naming the acting
      operator. The Clark--Wilson property~\cite{clark1987integrity} is the target,
      not the present state (\secref{sec:limitations}).
\end{itemize}

The honest caveat is that indirection is not free. A live target is a piece of
mutable global state that determines production behaviour, so its integrity becomes
critical. \sys's answer is that a rotation on the governed path is scope-checked,
audited, and checkpointed --- but that answer does not yet cover every route that
can perform one, which makes the trade sharper than it first appears.
\secref{sec:limitations} returns to it.

\subsection{R3 in practice: evidence bound to versions}

A score in a report tells you a number. A score attached to a version tells you
what you knew when you shipped. \sys binds per-dimension scores, the corpus
identity and version, the diagnosis, the selected optimizer, and the gate decision
to the candidate, and mints a provenance anchor at release carrying that bundle.

The mechanism that keeps the corpus honest over time is trace-to-example capture: a
verified production failure becomes durable test data, so the evidence base tracks
what production actually does rather than what it did when the suite was written.
This closes a specific drift problem that offline harnesses cannot close on their
own, because they have no access to the production trace stream.

\subsection{R4 in practice: one enforced gate, one human}

\secref{sec:components} defined the gate taxonomy; here we state what it buys, and
what it does not. The gate is enforced \emph{within the refinement state machine}:
a candidate that fails it never reaches review. It is not a system-wide interlock,
because the direct operator promotion endpoint can carry an advisory verdict and
an attributed override without entering that machine (\secref{sec:limitations}).
Authoring itself is non-live. Within the refinement path, the review
queue contains only candidates that passed their own evaluation. An operator's attention is therefore
spent on judgement rather than on triage --- which matters, because operator
attention is the binding constraint on this path (\secref{sec:arch}). Because the
per-version verdict is advisory, an urgent release is never blocked by a stale
artifact, and the override path that would otherwise erode the enforced gate never
needs to exist.

\subsection{Automation reach, stated as a boundary}

\sys automates diagnosis, candidate generation, evaluation, gating, evidence
assembly, and the mechanics of release. It does not automate the decision. The
reach is therefore: from a flagged trace to a scored, gated, reviewable candidate
without human involvement, and from a reviewed candidate to an audited release
without human involvement. The human sits at exactly one point in between, and at
one point before (verification).

We want to be precise that this is a design position rather than an unfinished
state. An auto-apply threshold would be straightforward to implement; D4 in
\tabref{tab:decisions} records why we consider it a mistake at the current state of
the art, and \secref{sec:future} records the conditions under which we would revisit
it.

\subsection{Cross-family capabilities}

\para{Knowledge bases.} Versioned retrieval corpora with chunking, embeddings,
hybrid retrieval combining lexical and dense scoring~\cite{robertson2009bm25,
lewis2020rag}, cross-encoder reranking~\cite{nogueira2019bertrerank}, and graph
extraction over Apache~AGE~\cite{apacheage}. The governance-relevant property is
that a build is an immutable version behind an active pointer, so ``which corpus
answered that question'' is a resolvable query rather than an archaeological one.

\para{Workflows.} A visual editor over a typed IR with \statNodeTypes{} node types,
queued runs, runtime approvals, checkpointing, and publication as an HTTP service.
The same IR backs the in-server interpreter and an Agents SDK export, which is what
lets a workflow be governed here and executed elsewhere.

\para{Tools and MCP servers.} Tool definitions carry deterministic test suites
fenced to a revision, so a suite result is attributable to a specific tool version.
MCP servers are managed definitions with discovered tool inventories, audited edit
history, fail-closed connection and policy controls, and a production preflight
that distinguishes containment from an attested boundary
(\secref{sec:components}).

\para{Observability.} \mlflow{} tracing with multimodal attachments --- extraction
spans carry the source document, so a knowledge-graph or OCR pipeline's inputs are
inspectable in the trace viewer --- plus server-sent live events, opt-in archival of
aged traces to object storage, and SLO reconciliation that produces durable
incidents~\cite{beyer2016sre}.

\subsection{Capability coverage against comparable systems}

%% fig-capability.tex -- Figure: capability coverage against comparable systems.
%%
%% A coverage grid, not a scorecard. Each cell states how far a system carries a
%% capability, using four levels that are defined once in the caption. Levels are
%% encoded by fill AND by glyph so the figure survives greyscale printing and is
%% legible to a reader who cannot distinguish the hues.
%%
%% Honesty constraint: this figure claims *architectural coverage* -- whether the
%% capability exists as a first-class, wired concern -- and claims nothing about
%% quality, performance, or maturity. Rows where CALIBER is not ahead are included
%% deliberately; a coverage grid with a winning column everywhere is advertising.

\begin{figure}[tbp]
\centering
\begin{cbwide}
\begin{tikzpicture}[x=1cm,y=1cm]

  \def\LW{5.40}    % capability label column width
  \def\CWD{1.44}   % per-system column width
  \def\RH{0.66}    % row height: two lines of \figB plus padding

  % level styles: fill + glyph, so hue is never the only channel
  \tikzset{
    lvNone/.style  = {absent, draw=cbRule},
    lvPart/.style  = {fill=cbMutedBg,    draw=cbMuted},
    lvAdj/.style   = {fill=cbSurfaceBg,  draw=cbSurface},
    lvFull/.style  = {fill=cbControlBg,  draw=cbControl},
  }
  % Offsets are in cm, not pt: inside a picture with x=1cm, pgfmath reads a bare
  % number as one x-unit, so writing "+1pt" here would shift a full centimetre.
  \newcommand{\lv}[4]{% col index, row index, style, glyph
    \node[#3, anchor=north west, minimum width=\CWD cm-3pt,
          minimum height=\RH cm-3pt, inner sep=0pt, line width=0.4pt,
          rounded corners=1pt, font=\figB, text=cbInk]
      at ({\LW+#1*\CWD+0.05}, {-#2*\RH-0.05}) {#4};}

  % ---- system headers ------------------------------------------------------
  \foreach [count=\i from 0] \n in {%
      {\sysbare}, {\mlflow{} 3\\registry}, {LangChain +\\LangGraph},
      {AutoGen \dt\\CrewAI}, {LangSmith \dt\\Langfuse},
      {DSPy}, {promptfoo \dt\\Evals}}
  {
    \node[anchor=south, font=\figB\bfseries, text width=\CWD cm-3pt,
          align=center, inner sep=1.5pt]
      at ({\LW+\i*\CWD+\CWD/2}, 0.08) {\n};
  }
  \draw[weak, draw=cbInk] (0,0) -- ({\LW+7*\CWD},0);

  % ---- rows ----------------------------------------------------------------
  \foreach [count=\r from 0] \cap/\gloss in {%
      {Typed registry, many families}/{beyond prompts: workflows, tools, corpora},
      {Live target indirection}/{an alias the client resolves at call time},
      {Enforced advancement gate}/{a failing evaluation stops the candidate},
      {Rollback from a record}/{the outgoing target is recorded, not inferred},
      {Provenance anchor per release}/{candidate and evidence bound to the version},
      {Human decision in the loop}/{an explicit apply action, not a threshold},
      {Automated candidate proposal}/{an optimizer proposes the next version},
      {Trace-to-example capture}/{production traces become test data},
      {LLM judges + human alignment}/{authored judges, agreement measured},
      {Governed tool execution}/{allowlists and a process boundary},
      {Governed MCP transport}/{fail-closed egress; containment or attestation},
      {RBAC on release}/{scope-checked live-target mutation},
      {Multi-agent orchestration}/{agents composed into a running system},
      {Managed cloud operation}/{someone else runs it for you}}
  {
    \node[anchor=north west, text width=\LW cm-7pt, font=\figB, align=left,
          inner sep=1.5pt] at (0,{-\r*\RH-0.03})
      {\textbf{\cap}\\{\itshape\gloss}};
    \draw[weak] (0,{-\r*\RH-\RH}) -- ({\LW+7*\CWD},{-\r*\RH-\RH});
  }

  % ---- the grid ------------------------------------------------------------
  % Columns: 0 CALIBER | 1 MLflow+LangSmith | 2 LangChain/Graph | 3 AutoGen/CrewAI
  %          4 Langfuse/Phoenix | 5 DSPy | 6 promptfoo/Evals
  % Levels:  Full = first-class and wired; Adj = present but partial or
  %          per-family; Part = achievable by assembling your own; None = absent.
  \foreach \c/\r/\s/\g in {%
    % typed asset registry -- MLflow registers prompts AND models; LangSmith and
    % Langfuse register prompts and datasets. Typed and first-class for those
    % types; the CALIBER claim is breadth of family, which the caption states.
    0/0/lvFull/{$\bullet$}, 1/0/lvFull/{$\bullet$}, 2/0/lvPart/{$-$},
    3/0/lvNone/{}, 4/0/lvFull/{$\bullet$}, 5/0/lvPart/{$-$}, 6/0/lvNone/{},
    % live target indirection -- MLflow aliases, LangSmith `production' tag,
    % Langfuse labels, all resolved client-side at call time. Not a differentiator.
    0/1/lvFull/{$\bullet$}, 1/1/lvFull/{$\bullet$}, 2/1/lvNone/{},
    3/1/lvNone/{}, 4/1/lvFull/{$\bullet$}, 5/1/lvNone/{}, 6/1/lvNone/{},
    % enforced advancement gate -- no documented mechanism makes a promotion fail
    % on a score; all three evaluate, none gate. CI can be assembled around them.
    0/2/lvFull/{$\bullet$}, 1/2/lvPart/{$-$}, 2/2/lvNone/{},
    3/2/lvNone/{}, 4/2/lvPart/{$-$}, 5/2/lvAdj/{$\circ$}, 6/2/lvAdj/{$\circ$},
    % audited release + rollback -- all three roll back by re-pointing; CALIBER
    % adds a recorded checkpoint, and only on the governed path.
    0/3/lvAdj/{$\circ$}, 1/3/lvAdj/{$\circ$}, 2/3/lvNone/{},
    3/3/lvNone/{}, 4/3/lvAdj/{$\circ$}, 5/3/lvNone/{}, 6/3/lvNone/{},
    % provenance anchor -- evidence linked to a version, but not bound as a
    % precondition of the promotion.
    0/4/lvFull/{$\bullet$}, 1/4/lvAdj/{$\circ$}, 2/4/lvNone/{},
    3/4/lvNone/{}, 4/4/lvAdj/{$\circ$}, 5/4/lvNone/{}, 6/4/lvNone/{},
    % human decision in the loop -- promotion is a deliberate human action in all
    % three; what differs is whether it is a step in a modelled state machine.
    0/5/lvFull/{$\bullet$}, 1/5/lvAdj/{$\circ$}, 2/5/lvPart/{$-$},
    3/5/lvPart/{$-$}, 4/5/lvAdj/{$\circ$}, 5/5/lvNone/{}, 6/5/lvNone/{},
    % automated candidate proposal -- MLflow ships optimize_prompts() with GEPA
    % and metaprompting. This row was wrong in an earlier draft.
    0/6/lvFull/{$\bullet$}, 1/6/lvFull/{$\bullet$}, 2/6/lvNone/{},
    3/6/lvNone/{}, 4/6/lvNone/{}, 5/6/lvFull/{$\bullet$}, 6/6/lvNone/{},
    % trace-to-example capture
    0/7/lvFull/{$\bullet$}, 1/7/lvAdj/{$\circ$}, 2/7/lvPart/{$-$},
    3/7/lvNone/{}, 4/7/lvAdj/{$\circ$}, 5/7/lvPart/{$-$}, 6/7/lvPart/{$-$},
    % LLM judges + human alignment
    0/8/lvFull/{$\bullet$}, 1/8/lvAdj/{$\circ$}, 2/8/lvPart/{$-$},
    3/8/lvNone/{}, 4/8/lvAdj/{$\circ$}, 5/8/lvPart/{$-$}, 6/8/lvAdj/{$\circ$},
    % governed tool execution
    0/9/lvAdj/{$\circ$}, 1/9/lvNone/{}, 2/9/lvPart/{$-$},
    3/9/lvPart/{$-$}, 4/9/lvNone/{}, 5/9/lvNone/{}, 6/9/lvNone/{},
    % governed MCP transport
    0/10/lvFull/{$\bullet$}, 1/10/lvNone/{}, 2/10/lvPart/{$-$},
    3/10/lvPart/{$-$}, 4/10/lvNone/{}, 5/10/lvNone/{}, 6/10/lvNone/{},
    % RBAC on release -- LangSmith owners-only tagging, Langfuse protected
    % labels; comparable to CALIBER's scope check, and none is Clark--Wilson.
    0/11/lvAdj/{$\circ$}, 1/11/lvAdj/{$\circ$}, 2/11/lvNone/{},
    3/11/lvNone/{}, 4/11/lvAdj/{$\circ$}, 5/11/lvNone/{}, 6/11/lvNone/{},
    % multi-agent orchestration -- CALIBER is deliberately not the leader here
    0/12/lvPart/{$-$}, 1/12/lvNone/{}, 2/12/lvFull/{$\bullet$},
    3/12/lvFull/{$\bullet$}, 4/12/lvNone/{}, 5/12/lvAdj/{$\circ$}, 6/12/lvNone/{},
    % managed cloud operation -- likewise
    % MLflow itself is a project rather than a service; managed operation comes
    % from a vendor distribution, so this is partial for the column as labelled.
    0/13/lvNone/{}, 1/13/lvAdj/{$\circ$}, 2/13/lvFull/{$\bullet$},
    3/13/lvPart/{$-$}, 4/13/lvFull/{$\bullet$}, 5/13/lvNone/{}, 6/13/lvPart/{$-$}}
  { \lv{\c}{\r}{\s}{\g} }

  % ---- column rules and the CALIBER column emphasis ------------------------
  \foreach \c in {0,...,7}
    { \draw[weak] ({\LW+\c*\CWD},0) -- ({\LW+\c*\CWD},{-14*\RH}); }
  \draw[draw=cbControl, line width=0.7pt, rounded corners=1.5pt]
    ({\LW-0.05},0.04) rectangle ({\LW+\CWD+0.05},{-14*\RH-0.04});

  % ---- legend --------------------------------------------------------------
  \node[anchor=north west, font=\figB, align=left, inner sep=0pt]
    at (0,{-14*\RH-0.26})
    {\tikz[baseline=-0.55ex]{\node[lvFull, minimum width=10pt, minimum height=7pt,
       inner sep=0pt, font=\figB]{$\bullet$};}~first-class and wired\qquad
     \tikz[baseline=-0.55ex]{\node[lvAdj, minimum width=10pt, minimum height=7pt,
       inner sep=0pt, font=\figB]{$\circ$};}~present but partial or per-family\qquad
     \tikz[baseline=-0.55ex]{\node[lvPart, minimum width=10pt, minimum height=7pt,
       inner sep=0pt, font=\figB]{$-$};}~achievable by assembling your
       own\qquad
     \tikz[baseline=-0.55ex]{\node[lvNone, minimum width=8pt, minimum height=6pt,
       inner sep=0pt]{};}~absent};

\end{tikzpicture}
\end{cbwide}
\caption{Capability coverage against comparable systems, \textbf{as of
2026-08-03}. This is a coverage grid, not a scorecard: a cell records whether a
capability exists as a first-class, wired concern, and says nothing about quality,
performance, or maturity. Levels are encoded by fill \emph{and} by glyph so the
figure reads in greyscale. Claims about the named products are sourced to dated
documentation and compared row by row in
\tabref{tab:platforms}~\cite{mlflowpromptregistry,mlflowoptimize,
langsmithprompts,langfuseprompts}; these are live products and the grid will age.
Note in particular that \sys shares --- rather than leads --- the pointer
indirection, rollback, and candidate-proposal rows: an earlier version of this
figure marked those against \mlflow{}, LangSmith, and Langfuse in error, and the
corrected grid gives \sys a much thinner lead. Where \sys is behind, the grid
also says so: it is deliberately not a multi-agent orchestration framework, and it
has no managed cloud offering, which for many teams is the deciding factor
(\secref{sec:compare}, \secref{sec:limitations}). A grid whose leftmost column
won every row would be advertising rather than evidence.}
\label{fig:capability}
\end{figure}


\figref{fig:capability} places these capabilities against the systems a reader is
most likely to be choosing between. It is a coverage grid rather than a scorecard: a
cell records whether a capability exists as a first-class, wired concern, and says
nothing about quality, performance, or maturity. Levels are encoded by fill
\emph{and} by glyph so the figure reads in greyscale.

The two rows where \sys is behind are included deliberately and are not
incidental. \sys is not a multi-agent orchestration framework and does not intend
to become one. It has no managed cloud offering, which for many teams is the
deciding factor regardless of architecture. A coverage grid whose leftmost column
won every row would be advertising rather than evidence, and a reader is entitled to
treat such a grid with suspicion.
